% !TeX spellcheck = en_US

\section{Data Used} \label{sec:dat}

This section describes the origin of the astronomical data used, the computational infrastructure implemented for their serialization, the systematic image quality auditing process, and the filtering strategy employed to construct the final clean catalog used for training our multimodal model.

\subsection{Initial Data Sample}

The initial dataset for this thesis is constructed from the intersection of two of the most significant cosmological surveys of the decade: the \textit{Euclid} space mission of the \gls{ESA} (DR1) and the \textit{\gls{DESI}} spectroscopic survey (DESI DR1, \autocite{ar:desi-early-release,ar:desi-dr1-2026}).

\subsubsection{Individual Selection Criteria}

The following quality filters were applied in parallel to the raw catalogs of both surveys:

\begin{enumerate}
\item \textbf{Quality filters on \textit{\gls{DESI}} DR1:}
    \begin{itemize}
    \item \texttt{ZCAT\_PRIMARY == True}: Ensures the selection of the best available spectrum for each unique physical object, eliminating redundant observations.
    \item \texttt{ZWARN == 0} or \texttt{ZWARN == 4}: Excludes sources with detected issues in the redshift ($z$) fitting, except those with minor non-critical spectral type warnings.
    \item \texttt{COADD\_FIBERSTATUS == 0}: Guarantees that the instrumental optical fiber did not present mechanical or readout anomalies during the exposure of the co-added spectrum.
    \item \texttt{SPECTYPE != 'STAR'}: Restricts the sample to extragalactic sources (galaxies and quasars), excluding Milky Way stars.
    \item \texttt{Z > 0}: Limits the sample to sources with positive cosmological redshifts ($z$).
    \end{itemize}

\item \textbf{Quality filters on Euclid DR1:}
    \begin{itemize}
    \item \texttt{VIS\_DET == 1}: Requires unambiguous detection of the source in the reference image of the visible (VIS) channel of \textit{Euclid} \autocite{ar:cropper-2024}.
    \item \texttt{spurious\_prob < 0.2}: Filters spurious detections produced by noise or local optical artifacts with a strict probability threshold.
    \item \texttt{flux\_detection\_total > 0.0}: Excludes objects with null or negative measured flux in the detection band.
    \end{itemize}
\end{enumerate}

\subsubsection{Symmetric Positional Crossmatch}

Once the quality criteria were applied to the individual catalogs, a symmetric best-match positional crossmatch was performed between the optical coordinates of \textit{Euclid} and those corresponding to the robotic fiber of \textit{\gls{DESI}}. A tolerance radius of $0.5''$ was established to avoid spurious or multiple-source matches in high-density fields.

The crossmatching process yielded the following results:

\begin{itemize}
\item \textbf{Wide-field sample (WIDE):} $206\,690$ physical matches.
\item \textbf{Deep-field sample (DEEP):} $46\,880$ physical matches.
\end{itemize}

Sources associated with \textit{Gaia} mission identifiers were not excluded during the crossmatch. This decision was made to preserve bright galaxies at low redshift ($z$) that automated stellar classification algorithms commonly misclassify as stars \autocite{ar:ecker-2026}. The final exclusion of Milky Way stars was entrusted to the \textit{\gls{DESI}} spectral classifier (\texttt{SPECTYPE}).

\subsubsection{Duplicate Resolution via Quality Hierarchy}

Due to the geographic overlap between the wide-field and deep-field surveys, $820$ unique objects (defined by their \texttt{TARGETID} identifier in \textit{\gls{DESI}}) were identified as being simultaneously present in both samples. To resolve these duplicates and ensure that the final catalog consisted solely of independent physical sources, a hierarchical selection scheme based on photometric quality and data provenance was implemented:

\begin{enumerate}
\item \textbf{Evaluation of \texttt{det\_quality\_flag} (Photometry):} The \textit{Euclid} detection quality indicator was prioritized in ascending order. Lower values indicate less degradation from defective pixels or sky background gradients. This criterion resolved $128$ duplicates ($15.6\%$).
\item \textbf{Evaluation of \texttt{spurious\_flag}:} In case of a tie in the previous step, the detection with the lowest spurious probability was prioritized (this produced no differences in the analyzed sample, as all objects presented the same value for this indicator).
\item \textbf{Evaluation by survey rank (\texttt{survey\_rank}):} If the tie persisted, data from the deep sample (DEEP, rank = 0) were prioritized over the wide-field sample (WIDE, rank = 1) due to the longer integration time and superior \textit{\gls{SNR}} of its observations. This step resolved the remaining $692$ duplicates ($84.4\%$).
\end{enumerate}

After deduplication, the initial combined sample was consolidated into $252\,750$ unique galaxies, of which $81.5\%$ ($205\,870$) originate from the WIDE sample, $18.2\%$ ($46\,060$) from the DEEP sample, and the remaining $0.3\%$ ($820$) from the overlap resolved via the quality hierarchy.

For each object, image cutouts were generated centered on their astronomical coordinates, with a fixed dimension of $224 \times 224$ pixels. These cutouts cover four \textit{Euclid} instrumental bands: the visible optical channel (VIS, $I_E$ band) of high spatial resolution, and the three near-infrared NISP channel bands ($Y_E$, $J_E$, and $H_E$ bands).

Since the nominal pixel scale of the VIS instrument is exactly $0.1$ arcseconds per pixel ($0.1''/\text{pixel}$) \autocite{ar:cropper-2024}, each $224 \times 224$ pixel cutout spans an effective field of view of $22.4'' \times 22.4''$. For the near-infrared NISP channel, whose native observations have a pixel scale of $0.3''/\text{pixel}$ \autocite{ar:jahnke-2025}, the images were bicubically interpolated and resampled onto the visible channel grid ($0.1''/\text{pixel}$). This ensures that all four bands of each cutout share the same matrix dimensions ($224 \times 224$ pixels) and precisely cover the same spatial extent of the sky ($22.4'' \times 22.4''$).

\subsection{Image Anomaly Filtering (Euclid)}

Despite the stringent restrictions applied during the crossmatching stage, massive astronomical catalogs inevitably contain undocumented instrumental anomalies. To identify and remove potential optical or readout corruptions in the \textit{Euclid} VIS and NISP image cutouts, an astronomical flux analysis module was designed and implemented.

This module audited the pixels of the cutouts for the $252\,744$ sources indexed in the \textit{Hugging Face Datasets} pretraining and test partitions. The algorithm classified the cutouts into four categories of physical integrity:

\begin{enumerate}
\item \textbf{Operational (correct):} Images with continuous and realistic flux values across all pixels in all four optical and infrared bands.
\item \textbf{Total null-flux anomalies (\textit{all-black}):} Cutouts whose pixels have strictly null values (zero or constant empty sky background values) across all four bands simultaneously. This indicates failures in the original mosaic cropping or areas without coverage at the detector edges.
\item \textbf{Partial null-flux anomalies (\textit{part-black}):} Cutouts displaying null pixel values in at least one instrumental band while retaining physical signal in the others (for example, a galaxy visible in VIS but with its infrared cutout in the $H_E$ band completely black).
\item \textbf{Non-numerical value corrupted samples (\textit{NaN-corrupted}):} Images containing undefined or infinite values (\texttt{NaN}/\texttt{Inf}) that destabilize gradient computation during backpropagation in neural network pretraining.
\end{enumerate}

\subsubsection{Audit Results}

The resulting integrity analysis log yielded the following statistics on the sample of $252\,744$ objects:

\begin{itemize}
\item Operational images (correct): $233\,944$ ($92.56\%$).
\item Total null-flux anomalies: $12\,486$ ($4.94\%$).
\item Partial null-flux anomalies: $6\,314$ ($2.50\%$).
\item Non-numerical value corrupted samples (\texttt{NaN}): $0$ ($0.00\%$).
\end{itemize}

The audit identified a cumulative total of $18\,800$ images with null-flux anomalies. The identifiers of these objects were recorded in an exclusion list for subsequent removal.

The conditional probability analysis of these anomalies revealed that there was no clear mathematical correlation between the default quality indicators of the \textit{Euclid} database (such as \texttt{det\_quality\_flag} or \texttt{spurious\_flag}) and the occurrence of null-flux cutouts. For example, the probability that an image was completely black given that it had been photometrically classified as free of photometric anomalies (\texttt{det\_quality\_flag == 0}) was $4.99\%$ ($8\,389$ out of $168\,127$), while the conditional probability for objects free of spurious detections (\texttt{spurious\_flag == 0}) was $4.94\%$. This demonstrates that the catalog indicators do not capture the physical absence of data in the locally cropped images, justifying the need for a pixel-level audit to purge the dataset.

\subsection{Final Pretraining Catalog Filtering}

To consolidate the final catalog and ensure physically consistent pretraining, a metadata filtering pipeline was executed on the combined FITS-format catalog of $252\,750$ records.

\paragraph{Phase 1: Adaptive Spectral SNR Criterion (DESI)} To eliminate noisy spectra that degrade loss optimization, a redshift-adaptive filter was applied to \textit{\gls{DESI}} data. For low-$z$ objects, key emission lines (e.g., $\text{H}\alpha$ and $\left[\text{N\,II}\right]$) fall in the red camera ($R$). For $z \ge 0.15$, cosmological redshift shifts them to the near-infrared $Z$ camera. Mathematically, the spectrum acceptance condition is:

\begin{equation}
\text{SNR Filter} = \left( \left(z < 0.15\right) \wedge \left(\text{SNR\_SPEC\_R} > 3.0\right) \right) \vee \left( \left(z \ge 0.15\right) \wedge \left(\text{SNR\_SPEC\_Z} > 3.0\right) \right)
\label{eq:filtro_snr}
\end{equation}

where \texttt{SNR\_SPEC\_R} and \texttt{SNR\_SPEC\_Z} are the average \textit{\gls{SNR}} in the $R$ and $Z$ cameras. This filter accepted $145\,660$ objects ($57.6\%$) and rejected $107\,090$ objects ($42.4\%$) from the initial catalog.

\paragraph{Phase 2: Corrupt Image Exclusion (Euclid)} The $145\,660$ objects with scientific-quality spectroscopy were crossmatched against the $18\,800$ corrupt \textit{Euclid} images identified during the pixel audit. This discarded $7\,999$ objects that had passed the spectral \gls{SNR} filter but possessed defective image cutouts.


\begin{figure}[ht]
	\centering
	\includegraphics[width=\textwidth]{target_diagnostic_39627829519845965}
	\caption[Example of one of the objects in the dataset]{Visualization of one of the objects that make up the dataset. In the top center, the false-color image from \textit{Euclid} is shown. The color processing follows the methodology proposed by \textcite{ar:lupton-2004} for the \textit{\gls{SDSS}} but adapted to the four channels of \textit{Euclid}. On the top sides, different galaxy properties are displayed. At the bottom, divided into two channels, the spectrum obtained by the \textit{\gls{DESI}} spectrographs is shown. Some of the spectral lines used as references throughout the project are plotted.}
	\label{fi:galaxy_bimodal_sample}
\end{figure}


\paragraph{Final Clean Catalog Consolidation}

After sequentially applying the \textit{\gls{DESI}} spectral \gls{SNR} filter and the \textit{Euclid} image quality rejection, the final clean catalog was consolidated to $\mathbf{137\,661}$ unique galaxy image-spectrum pairs. This final catalog represents a net retention of $54.5\%$ relative to the initial overlap sample and was consolidated into a unified binary FITS file.

This clean catalog of $137\,661$ objects constitutes the optimal statistical foundation for pretraining the multimodal foundation model \textit{AstroPT}. By combining artifact-free images with high-\gls{SNR} spectra, it is ensured that the unified latent space is structured under real physical principles, minimizing the memorization of instrumental noise or data reduction systematics.

It is important to note that in diagnostic analyses performed after the consolidation of this sample, a minor subset of \textit{Euclid} images was detected that exhibited morphological corruptions due to a physical misalignment in the registration of the filters (VIS versus NISP). These sources retain non-null pixel values across all their bands and consequently evade automatic detection by the flux analysis algorithm, which only searches for null or black values. However, it was intentionally decided not to purge these objects from the final catalog. The rationale behind this decision lies in the desire to train the foundation model \textit{AstroPT} under realistic and non-idealized observational conditions, forcing it to latently isolate chromatic misalignments and instrumental domain imperfections as intrinsic anomalies or artifacts, thereby dispensing with destructive preprocessing or excessively restrictive manual filtering.

\subsubsection{Origin and Determination of Reference Values}

The establishment of reference values (\textit{ground truth}) for each of the physical parameters in our consolidated catalog is based on the synergistic integration of the official \glspl{VAC} and data releases from both surveys.

\paragraph{Structural and Photometric Data (\textit{Euclid} DR1):}                                                                                                                      
	The analytical geometric parameters ($R_{\text{eff}}$, $n_{\text{Sersic}}$, and $q$) are directly queried from the official photometric catalog of the \textit{Euclid} DR1. These structural parameters were computed by the \textit{Euclid} Consortium pipeline through two-dimensional analytical Sérsic brightness profile fitting on the visible VIS band images, utilizing morphological extraction algorithms such as \textit{SExtractor} \autocite{ar:bertin-1996} or \textit{GALFIT} \autocite{ar:peng-2002}. 
	
	The non-parametric morphology metrics (Gini and smoothness) are also extracted from the DR1 catalogs, derived from automated morphometry codes. The qualitative morphological tracers (e.g., the probability of possessing spiral arms, \texttt{has\_spiral\_arms\_yes}) are sourced from the Galaxy Zoo DESI \gls{VAC} \autocite{ar:walmsley-2023}, where convolutional networks assigned probabilities to each source.
	
\paragraph{Spectroscopic and Integrated Data (\textit{DESI} DR1 VACs):} 
	The integrated physical and kinematic properties are strictly retrieved from the official \glspl{VAC} of the \textit{DESI} First Data Release (DR1) \autocite{ar:desi-early-release,ar:desi-dr1-2026,ar:siudek-2024}. 
	
	Specifically, the stellar physical properties ($\log M_*$ and $\log \text{\gls{SFR}}$) are extracted from the physical properties \gls{VAC} \autocite{ar:siudek-2024}. In their pipeline, these parameters are estimated via the \textit{CIGALE} \autocite{ar:boquien-2019} code, which performs a spectrophotometric synthesis fit combining the high-resolution spectrum of \textit{DESI} with broadband optical photometry from the Legacy Surveys \autocite{ar:dey-2019} supplemented by mid-infrared photometry from \textit{\gls{WISE}}\autocite{ar:wright-2010}.
	
	The fluxes and equivalent widths of the spectroscopic emission lines ($F(\text{H}\alpha)$, $F(\text{H}\beta)$, etc.) and the gas kinematics ($\sigma_{[\text{O\,III}]}$) are sourced from the \textit{FastSpecFit} \gls{VAC} \autocite{ar:moustakas-2023}. The \textit{FastSpecFit} automated pipeline models the stellar continuum using high-resolution templates and applies simultaneous Gaussian fits to the emission lines, providing robust measurements that correct for internal dust attenuation and stellar atmospheric absorption.

\subsection{Reference Physical Parameters and Ground Truth} \label{sec:dat:params}

To quantitatively evaluate the physical information encoded in the latent representations obtained during multimodal pretraining, a set of $16$ fundamental astrophysical parameters was selected. These parameters serve as target variables in the regression tests determining the quality of the models described in \cref{sec:dev}. It is critical to emphasize that these parameters are exclusively derived from the official data releases and \glspl{VAC} of the \textit{Euclid} and \textit{\gls{DESI}} pipelines; they are not computed locally, but act as our absolute empirical ground truth. A brief description of each parameter follows.

\subsubsection{Structural and Photometric Data}

\begin{itemize}
	\item \textbf{Effective elliptical Sérsic radius ($R_{\text{eff}}$, column \texttt{sersic\_sersic\_vis\_radius}):} The radius measured along the semi-major axis that encloses $50\%$ of the projected luminosity of the galaxy, measured in arcseconds. It represents the intrinsic physical size and angular scale of the source on the sky.
	\item \textbf{Sérsic index ($n_{\text{Sersic}}$, column \texttt{sersic\_sersic\_vis\_index}):} A structural parameter that determines the surface brightness profile \autocite{ar:sersic-1968}. Values near $n_{\text{Sersic}} \approx 1$ describe exponential disks characteristic of late-type (spiral) galaxies, while values near $n_{\text{Sersic}} \approx 4$ define de Vaucouleurs profiles \autocite{ar:vaucouleurs-1959} typical of compact bulges or early-type elliptical galaxies.
	\item \textbf{Axis ratio ($q = b/a$, column \texttt{sersic\_sersic\_vis\_axis\_ratio}):} The ratio between the semi-minor axis ($b$) and the semi-major axis ($a$) of the elliptical fit. It measures the degree of intrinsic ellipticity and the spatial inclination of the galactic disk with respect to the observer.
	\item \textbf{Gini coefficient ($G$, column \texttt{gini}):} A non-parametric statistical measure of luminous inequality. It quantifies the degree of uniformity or homogeneity in the distribution of luminous flux among the galaxy's pixels. Values close to $1$ indicate strong luminous concentration in a few pixels (for example, compact bulges or active galactic nuclei), while low values reflect dispersed and uniform brightness.
	\item \textbf{Photometric smoothness or clumpiness (Smoothness, $S$, column \texttt{smoothness}):} A non-parametric metric that measures the fraction of light associated with small-scale discrete structures. It is calculated by subtracting a low-pass filtered smoothed version from the original image, revealing young stellar clusters and active peripheral H II regions.
	\item \textbf{Spiral arms fraction (column \texttt{has\_spiral\_arms\_yes}):} A non-parametric fraction obtained from convolutional classifiers trained on galaxy morphology, measuring the conditional probability that the galaxy exhibits active and structured spiral arms in the visual cutout.
	\item \textbf{Detection photometric flux (\texttt{flux\_detection\_total}):} The integrated apparent brightness measured in the \textit{Euclid} VIS band, expressed in microjanskys ($\mu\text{Jy}$), representing the total detection flux.
\end{itemize}

\subsubsection{Spectroscopic and Integrated Parameters}

\begin{itemize}
	\item \textbf{Cosmological redshift ($z$, column \texttt{Z}):} The spectroscopic shift of atomic lines relative to their rest-frame wavelengths. It provides a clean, shortcut-free measurement of the recession velocity due to the Hubble flow and the distance of the galaxy.
	
	\item \textbf{Integrated stellar mass ($\log M_*$, column \texttt{LOGMSTAR}):} The total physical mass accumulated in the form of stable stars in the galaxy, expressed on a logarithmic scale in solar masses ($\log M_\odot$).
	
	\item \textbf{Star Formation Rate (\textit{log \gls{SFR}}, column \texttt{LOGSFR}):} The rate at which the galaxy transforms its cold molecular gas into new stars, expressed in logarithm of solar masses per year ($\log M_\odot/\text{yr}$).
	
	
	\item \textbf{$\text{H}\alpha$ line flux ($F(\text{H}\alpha)$, column \texttt{HALPHA\_FLUX}):} Integrated flux of the Balmer hydrogen recombination line $\text{H}\alpha$ at $6563\,\text{\AA{}}$, expressed in standard units of $10^{-17}\,\text{erg}\,\text{s}^{-1}\,\text{cm}^{-2}$.
	
	\item \textbf{$\text{H}\alpha$ equivalent width ($EW(\text{H}\alpha)$, column \texttt{HALPHA\_EW}):} The ratio of the integrated emission line flux of $\text{H}_\alpha$ normalized by the flux density of the surrounding stellar continuum, expressed in Angstroms ($\text{\AA{}}$). It is indicative of the specific star formation rate.
	
	\item \textbf{$\text{H}\beta$ line flux ($F(\text{H}\beta)$, column \texttt{HBETA\_FLUX}):} Integrated flux of the Balmer recombination line $\text{H}_\beta$ at $4861\,\text{\AA{}}$, expressed in standard units of $10^{-17}\,\text{erg}\,\text{s}^{-1}\,\text{cm}^{-2}$. It is crucial for correcting internal interstellar absorption and dust reddening through the Balmer decrement.
	
	\item \textbf{$[\text{N\,II}]\,\lambda 6584$ line flux ($F([\text{N\,II}])$, column \texttt{NII\_6584\_FLUX}):} Flux of the forbidden atomic transition of singly ionized nitrogen, expressed in standard units of $10^{-17}\,\text{erg}\,\text{s}^{-1}\,\text{cm}^{-2}$. It is sensitive to the metallicity of the ionized gas and to non-stellar ionization mechanisms, such as those produced by \textit{\glspl{AGN}}.
	
	\item \textbf{$[\text{O\,III}]\,\lambda 5007$ line flux ($F([\text{O\,III}])$, column \texttt{OIII\_5007\_FLUX}):} Flux of the doubly ionized oxygen line, expressed in standard units of $10^{-17}\,\text{erg}\,\text{s}^{-1}\,\text{cm}^{-2}$. It is an indicator of the presence of hard ionizing radiation from very young stellar clusters or from accretion disks of supermassive black holes.
	
	\item \textbf{$[\text{O\,III}]\,\lambda 5007$ velocity dispersion ($\sigma_{[\text{O\,III}]}$, column \texttt{OIII\_5007\_SIGMA}):} Doppler spectral width of the $[\text{O\,III}]$ line, expressed in kilometers per second ($\text{km}\,\text{s}^{-1}$), which measures the turbulent kinematics and velocity dispersion of the nuclear ionized gas induced by outflows or by the central gravitational potential of the galaxy.
\end{itemize}

\subsection{Data Storage Format}

The efficient storage and processing of hundreds of thousands of high-resolution spectral pairs and multi-band image cutouts represents a significant challenge in \textit{\gls{HPC}} environments. Traditional astronomical file storage formats (such as individual binary FITS tables) introduce bottlenecks in the input/output (I/O) data transfer rate due to the latency involved in performing concurrent random reads of millions of small files on disk.

To mitigate this I/O limitation, the dataset was serialized in columnar \textit{Apache Arrow} format using the \textit{Hugging Face Datasets} library. This paradigm offers the following architectural advantages:

\begin{enumerate}
	\item \textbf{High-Speed Sequential Reading:} Groups images and spectra into large memory-mapped files, enabling continuous reads at speeds approaching the physical limit of SSD/NVMe storage.
	\item \textbf{Zero-Copy Serialization:} \textit{Apache Arrow} allows the byte arrays loaded into host RAM to be shared directly with the \textit{PyTorch} data loading subprocesses (dataloaders) without requiring memory duplication, reducing RAM consumption.
	\item \textbf{Efficient Random Access:} Enables shuffling of samples during pretraining on \glspl{GPU} instantaneously without suffering penalties from random disk access.
\end{enumerate}

\subsubsection{Dataset Partitioning}

The consolidated dataset in \textit{Apache Arrow} format was randomly divided into two independent partitions for model pretraining and validation, using a fixed pseudorandom seed to ensure experimental reproducibility:

\begin{itemize}
	\item Pretraining set (train split): $202\,196$ objects ($80\%$).
	\item Validation and test set (test split): $50\,548$ objects ($20\%$).
\end{itemize}

This partition totals $252\,744$ image-spectrum pairs that were indexed in the \textit{Hugging Face Datasets} columnar format.